\chapter{Що роблять нейрони уві сні}
\chaptermark{Сон}
\label{sec:sleep}

\section{Сон --- не вимкнення, а інший режим}

Інженер, побачивши вимкнений комп'ютер, скаже: він нічого не робить. Зі сплячим мозком так не виходить. Нейрони уві сні не мовчать: у глибокому сні кора видає імпульси, у швидкому сні --- майже стільки ж, скільки вдень. Змінюється не те, чи \emph{працюють} нейрони, а те, \emph{як} вони працюють. Сон --- це набір режимів машини, і перемикають їх ті самі широкомовні канали, що й удень.

Для кожного режиму можна виписати «стан регістрів» (таблиця~\ref{tab:sleepmodes}). Удень \nt{ACET}, \nt{NORA} і \nt{SERO} працюють, датчики підключено, м'язи слухаються. У повільному, глибокому сні всі три канали затихають, а вхід від органів чуття ослаблено~\cite{Hasselmo1999}: викид ацетилхоліну в корі падає, а \nt{NORA}- і \nt{SERO}-нейрони видають імпульси рідше, ніж удень~\cite{Marrosu1995,AstonJonesBloom1981,McGintyHarper1976}. У швидкому сні картина дивна: \nt{ACET} знову піднімається до денного рівня~\cite{Marrosu1995}, а \nt{NORA} і \nt{SERO} замовкають майже повністю~\cite{AstonJonesBloom1981,McGintyHarper1976,JacobsAzmitia1992}. М'язи тіла у швидкому сні вимкнено: \nt{ACET}-нейрони, що ними керують, сильно гальмуються через \nt{GABA} і \nt{GLYC}~\cite{BrooksPeever2012} --- тією самою забороною з розділу~\ref{sec:gaba}, тільки накладеною на весь вихід одразу.

\begin{table}[t]
\centering
\footnotesize
\setlength{\tabcolsep}{4pt}
\renewcommand{\arraystretch}{1.15}
\begin{tabular}{>{\raggedright}p{3.1cm}>{\raggedright}p{2.9cm}>{\raggedright}p{3.2cm}>{\raggedright\arraybackslash}p{3.4cm}}
\toprule
& неспання & повільний сон & швидкий сон \\
\midrule
\nt{ACET} (запис, розд.~\ref{sec:acet}) & високий & низький & високий \\
\nt{NORA} (підсилення, розд.~\ref{sec:nora}) & середній, спалахи & низький & майже мовчить \\
\nt{SERO} (розд.~\ref{sec:serocomputer}) & середній & низький & майже мовчить \\
вхід від датчиків & підключено & ослаблено & ослаблено \\
вихід на м'язи & підключено & ослаблено & вимкнено гальмуванням \\
активність кори & рівна & повільні гойдалки: робота --- тиша & рівна, як удень \\
\bottomrule
\end{tabular}
\caption{Стан регістрів машини в трьох режимах. Усі рядки виміряно; рівні \nt{ACET}, \nt{NORA} і \nt{SERO} --- прямими записами за стадіями сну~\cite{Marrosu1995,AstonJonesBloom1981,McGintyHarper1976}.}
\label{tab:sleepmodes}
\end{table}

\section{Коли спати: реле з гістерезисом}

Що вирішує, коли машині перемикатися в сон? Добре працює проста схема з двох частин~\cite{Borbely1982}. Перша частина --- \emph{тиск сну} $S$, який накопичується, поки ми не спимо, і скидається, поки спимо. Це конденсатор, що заряджається вдень і розряджається вночі:
\begin{empheq}[box=\keybox]{equation}
\frac{\dd S}{\dd t} \;=\; \frac{1-S}{\tau_r}\ \ \text{(неспання)},\qquad
\frac{\dd S}{\dd t} \;=\; -\frac{S}{\tau_d}\ \ \text{(сон)} .
\label{eq:sleeppressure}
\end{empheq}
Друга частина --- два пороги: машина засинає, коли $S$ доходить до верхнього порога $H(t)$, і прокидається, коли $S$ опускається до нижнього $L(t)$. Обидва пороги плавно гойдаються в такт добовому ритму з розділу~\ref{sec:chemoutput}. Вийшло реле з гістерезисом --- тригер Шмітта з підрозділу~\ref{sec:bistable}, --- а разом із конденсатором воно дає релаксаційний генератор, що підстроюється за фазою добовим ритмом~\eqref{eq:pll}.

\begin{figure}[t]
\centering
\begin{tikzpicture}[>=Stealth,x=0.16cm,y=4.2cm]
\foreach \a/\b in {0/4.3,21.2/29.3,45.4/53.4,69.5/72}{\fill[blue!8] (\a,0) rectangle (\b,1);}
\draw[->,gray!70!black] (0,0) -- (75,0) node[right,font=\small,black] {$t$, год};
\draw[->,gray!70!black] (0,0) -- (0,1.08) node[left,font=\small,black] {$S$};
\foreach \t in {0,24,48,72}{\draw[gray!60!black] (\t,-0.02) -- (\t,0.02); \node[below,font=\scriptsize] at (\t,0) {\t};}
\draw[thick,densely dashed,red!70!black] plot file {figures/sleep_H.dat};
\draw[thick,densely dashed,green!45!black] plot file {figures/sleep_L.dat};
\draw[very thick,blue!60!black] plot file {figures/sleep_S.dat};
\node[font=\scriptsize,red!70!black,anchor=west] at (73,0.72) {$H$: заснути};
\node[font=\scriptsize,green!45!black,anchor=west] at (73,0.24) {$L$: прокинутися};
\node[font=\scriptsize,blue!60!black] at (25.25,0.93) {сон};
\node[font=\scriptsize,blue!60!black] at (49.4,0.93) {сон};
\end{tikzpicture}
\caption{Тиск сну $S$~\eqref{eq:sleeppressure} за $\tau_r=18.2$~год, $\tau_d=4.2$~год. Удень $S$ заряджається до верхнього порога $H$, уночі (зафарбовано) розряджається до нижнього $L$; обидва пороги гойдаються разом із добовим ритмом. Машина сама приходить до режиму близько $16$ годин неспання і $8$ годин сну.}
\label{fig:twoprocess}
\end{figure}

У нашій моделі зі звичайними значеннями $\tau_r=18.2$~год і $\tau_d=4.2$~год машина сама виходить на $16.1$ години неспання і $8.0$ години сну (рис.~\ref{fig:twoprocess}). Модель пояснює й те, що здається дивним: після сорока годин без сну тиск доходить до $0.90$, але відновлювальний сон у моделі триває лише $8.8$ години, а не на добу довше за звичайний. Розряд експоненційний, і високий заряд спадає швидко; «борг» повертається не тривалістю, а глибиною сну.

Що фізично грає роль $S$? Сильний кандидат --- аденозин, «лічильник навантаження» з додатка~\ref{sec:othermediators}: його концентрація зростає за час неспання і падає уві сні~\cite{PorkkaHeiskanen1997,Dunwiddie2001}. Що тиск сну --- це саме аденозин і нічого більше, --- гіпотеза.

\section{Повільний сон: уся мережа як релаксаційний генератор}

У глибокому сні нейрони кори працюють почергово всією групою: рідше ніж раз на секунду вони разом умикаються на кілька сотень мілісекунд (\emph{стан роботи}) і разом замовкають (\emph{стан тиші}): частота цих гойдалок нижча за $1$~Гц, під наркозом найчастіше $0.3$--$0.4$~Гц~\cite{Steriade1993}. Ці повільні гойдалки ми можемо зібрати з деталей, які вже маємо. Візьмімо групу \nt{GLUT}-нейронів із сильним зворотним зв'язком на себе --- пам'ять із розділу~\ref{sec:glutcomputer}, яка має два стійкі стани, --- і додамо повільну втому, як у генераторі ритму розділу~\ref{sec:muscles}:
\begin{empheq}[box=\keybox]{equation}
\tau\,\frac{\dd r}{\dd t} = -r + F\bigl(w\,r + I - a\bigr),
\qquad
\tau_a\,\frac{\dd a}{\dd t} = -a + b\,r .
\label{eq:updown}
\end{empheq}
Група вмикається, працює й утомлюється; втома $a$ з'їдає верхній стан, і група падає в тишу; у тиші втома минає, нижній стан зникає, і група знову вмикається. Виходить той самий релаксаційний генератор, що й тиск сну, тільки приблизно у вісімдесят тисяч разів швидший.

\begin{figure}[t]
\centering
\begin{tikzpicture}[>=Stealth,x=2.9cm,y=1.0cm]
\begin{scope}[yshift=1.9cm]
\draw[->,gray!70!black] (0,0) -- (4.1,0);
\draw[->,gray!70!black] (0,0) -- (0,1.25) node[left,font=\small,black] {$r$};
\draw[thick,blue!60!black] plot file {figures/updown_sleep.dat};
\node[font=\scriptsize,anchor=west] at (4.15,0.5) {сон: $b=5$};
\end{scope}
\draw[->,gray!70!black] (0,0) -- (4.1,0) node[right,font=\small,black] {$t$, с};
\draw[->,gray!70!black] (0,0) -- (0,1.25) node[left,font=\small,black] {$r$};
\draw[thick,red!70!black] plot file {figures/updown_wake.dat};
\node[font=\scriptsize,anchor=west] at (4.15,0.5) {день: $b=1$};
\foreach \t in {0,1,2,3,4}{\draw[gray!60!black] (\t,-0.05) -- (\t,0.05); \node[below,font=\scriptsize] at (\t,0) {\t};}
\end{tikzpicture}
\caption{Та сама група~\eqref{eq:updown} у двох режимах: $w=5$, $I=2$, $\theta=2.5$, $\tau=10\ms$, $\tau_a=0.3$~с. Із сильною втомою ($b=5$, як у повільному сні) група гойдається між роботою й тишею з періодом $1.1$~с (значення моделі; у корі період більший за секунду, під наркозом --- близько $3$~с). Послабмо втому ($b=1$) --- так діє ацетилхолін --- і та сама група працює рівно.}
\label{fig:updown}
\end{figure}

А що переводить мережу з гойдалок у рівну денну роботу? Ацетилхолін пригнічує в нейронах повільні струми, що створюють утому~\cite{McCormick1992}. У нашій моделі за сильної втоми, $b=5$, група гойдається з періодом $1.1$~с --- швидше за виміряні гойдалки, але того самого порядку --- і працює близько $35\%$ часу, якщо усереднювати за цілим числом періодів; за слабкої, $b=1$, та сама група працює рівно (рис.~\ref{fig:updown}). Один регістр --- рівень \nt{ACET} --- перемикає генератор на підсилювач. Поверх повільних гойдалок уві сні йдуть і швидші спалахи ритму $7$--$15$ коливань на секунду; їх породжує петля \nt{GLUT}--\nt{GABA} між корою й проміжним релейним вузлом~\cite{SteriadeMcCormickSejnowski1993}, родичка ритму з розділу~\ref{sec:gaba}.

\section{Повтори: перемотування дня}

У розділі~\ref{sec:acet} ми бачили, що в повільному сні нейрони, які працювали вдень разом, знову спрацьовують разом і в тому самому порядку. Додамо одну інженерну деталь: повтори йдуть \emph{із перемотуванням}. Послідовність, яка вдень займала секунди, програється уві сні за десяті частки секунди~\cite{Nadasdy1999}: у гіпокампі приблизно вдвадцятеро швидше~\cite{LeeWilson2002}, у корі --- у шість-сім разів~\cite{Euston2007}. І програється не будь-коли: короткі спалахи повторів в одній ділянці підлаштовуються під гойдалки роботи й тиші та під швидкі спалахи ритму в корі. Якщо штучно посилити цю узгодженість, пам'ять поліпшується~\cite{Maingret2016} --- це вже причиновий зв'язок, а не збіг.

Навіщо машині перемотувати день? Інженерові знайома трудність, яку називають \emph{катастрофічним забуванням}: мережа, яка вчить нове підряд, одне за одним, псує старе. Рятує \emph{перемішування}: учити нове впереміш зі старим, невеликими порціями, багато разів. Гіпотеза про дві системи навчання~\cite{McClelland1995} полягає в тому, що швидка пам'ять записує день цілком, а уві сні потроху, впереміш і багаторазово програє його повільній корі. Це та сама пара «швидкий буфер --- повільне сховище», що в розділі~\ref{sec:acet}, і та сама пара швидкого й повільного учнів, що в розділі~\ref{sec:motorlearning}. Повтори та їхню узгодженість виміряно; що їхнє призначення --- перемішане навчання повільної пам'яті, --- добре обґрунтована гіпотеза.

\section{Перенормування ваг}

Удень правила Гебба з розділу~\ref{sec:dofa} здебільшого посилюють зв'язки. Якщо лише посилювати, ваги ростуть, а разом із ними ростуть витрати енергії й падає чутливість: нейрон, у якого всі входи сильні, перестає їх розрізняти. За гіпотезою \emph{синаптичного гомеостазу}~\cite{TononiCirelli2014}, сон потрібен, зокрема, для того, щоб повернути ваги до робочого рівня: усі зв'язки слабшають в одну й ту саму кількість разів, тож їхні \emph{відношення} --- вивчене --- зберігаються. Для інженера це нормування, як у правилі~\eqref{eq:oja}, тільки зроблене не на ходу, а в окремий час.

Прямі вимірювання цьому не суперечать: у мишей площа контакту на синапсах кори після сну приблизно на $18\%$ менша, ніж після неспання, зменшення пропорційне розміру, а найбільші й найстійкіші контакти майже не змінюються~\cite{deVivo2017}. Масштабування ваг виміряно; що це головне завдання сну, --- гіпотеза.

\section{Швидкий сон: машина, відключена від входів і виходів}

У швидкому сні кора працює майже як удень, але вхід від органів чуття ослаблено, а вихід на м'язи вимкнено гальмуванням. Для інженера це знайомий режим: \emph{випробування на стенді}. Схема працює на повну потужність, але її входи замкнено на внутрішній генератор, а виходи від'єднано від виконавчих механізмів, щоб нічого не зламати. Сновидіння, за однією з гіпотез, --- саме такий внутрішній прогін моделі світу, зібраної за день; що саме мережа при цьому робить і чи вчиться, невідомо. Виміряно тут лише те, що записано в таблиці~\ref{tab:sleepmodes}: який канал увімкнено, який мовчить, і що вихід заблоковано.

\section{Обслуговування і місцевий сон}

Довго вважалося, що уві сні міжклітинні щілини розширюються і мозок швидше промивається від продуктів обміну речовин~\cite{Xie2013}. Недавні досліди, що вимірювали промивання іншим способом, дістали протилежне: уві сні воно повільніше~\cite{Miao2024}. Питання нині відкрите, і ми його залишимо відкритим.

Надійніший інший факт: сон --- властивість місцевих мереж, а не всієї машини одразу. Якщо тварині довго не давати спати, то окремі ділянки її кори, поки тварина не спить і рухається, ненадовго провалюються в тишу, як у повільному сні, і в ці моменти тварина частіше помиляється~\cite{Vyazovskiy2011}. Кожна мережа втомлюється сама й сама перемикається; загальний режим виходить тоді, коли широкомовні канали переводять усі мережі одразу.

\begin{keyblock}
\begin{proposition}[Сон як набір режимів]
\label{prop:sleep}
Уві сні нейрони не вимикаються: широкомовні канали переводять машину в інші режими. Коли спати, вирішує реле з гістерезисом~\eqref{eq:sleeppressure}, підстроєне добовим ритмом. У повільному сні мережа стає релаксаційним генератором~\eqref{eq:updown} і програє день із перемотуванням; у швидкому --- працює, відключена від входів і виходів. Режими, повтори й масштабування ваг виміряно; їхнє призначення --- перемішане навчання й перенормування --- добре обґрунтовані гіпотези.
\end{proposition}
\end{keyblock}

\section{Підсумок}

\begin{table}[t]
\centering
\footnotesize
\setlength{\tabcolsep}{4pt}
\renewcommand{\arraystretch}{1.15}
\begin{tabular}{>{\raggedright}p{3.2cm}>{\raggedright}p{4.4cm}>{\raggedright\arraybackslash}p{5.2cm}}
\toprule
Що відбувається & Як & Що відомо \\
\midrule
зміна режимів & регістри \nt{ACET}, \nt{NORA}, \nt{SERO} (табл.~\ref{tab:sleepmodes}) & виміряно \\
коли спати & конденсатор $S$ і реле з гістерезисом~\eqref{eq:sleeppressure} & модель добре описує досліди; аденозин як $S$ --- гіпотеза \\
повільні гойдалки & група зі зворотним зв'язком і втомою~\eqref{eq:updown} & гойдалки виміряно; модель --- найпростіша схема \\
перемотування дня & повтори в $6$--$20$ разів швидші, узгоджені з гойдалками & виміряно; посилення узгодженості поліпшує пам'ять \\
навіщо повтори & перемішане навчання повільної пам'яті & добре обґрунтована гіпотеза \\
перенормування ваг & масштабування зв'язків уві сні & зменшення контактів виміряно; головна роль сну --- гіпотеза \\
швидкий сон & робота з відключеними входами й виходом & режим виміряно; призначення невідоме \\
промивання & розширення міжклітинних щілин & суперечливі результати \\
місцевий сон & окремі ділянки провалюються в тишу & виміряно \\
\bottomrule
\end{tabular}
\caption{Що роблять нейрони уві сні.}
\label{tab:sleep}
\end{table}

Таблиця~\ref{tab:sleep} підсумовує розділ. Сон виявився не паузою в роботі машини, а її обслуговуванням на ходу: перемикання режимів тими самими широкомовними каналами, генератор із тих самих деталей, що й ритм кроків, перемотування дня для повільної пам'яті й перенормування ваг. Удень машина пише, уночі переписує й упорядковує записане.

\section{Задачі}

\begin{exercise}[\lvlA]
\label{ex:20:1}
\emph{Реле без добового ритму.} Нехай пороги сталі: $H=0.67$, $L=0.17$ (середні значення порогів моделі рис.~\ref{fig:twoprocess}). Виведіть із~\eqref{eq:sleeppressure} тривалості неспання й сну та період генератора за $\tau_r=18.2$~год, $\tau_d=4.2$~год. Чому модель із гойдальними порогами все ж виходить на $16.1+8.0=24$~год? До якого приладу розділу~\ref{sec:chemoutput} належить цей ефект?
\end{exercise}

\begin{exercise}[\lvlA]
\label{ex:20:2}
\emph{Борг сну.} Сон, розпочатий за тиску $S_0$, триває $t_s=\tau_d\ln(S_0/L)$, $\tau_d=4.2$~год, $L$ --- нижній поріг під час пробудження. (а)~Після $40$~год без сну $S_0=0.90$, а сон тривав $8.8$~год. Яким був $L$? Скільки тривав би сон за звичайного $L\approx0.092$? (б)~За ніч площа синаптичних контактів зменшується на $18\%$. На скільки мають зрости ваги за день?
\end{exercise}

\begin{exercise}[\lvlB]
\label{ex:20:3}
\emph{Проєктування порогів.} Доберіть сталі пороги $H$ і $L$, за яких реле~\eqref{eq:sleeppressure} з $\tau_r=18.2$~год і $\tau_d=4.2$~год дає рівно $16$~год неспання і $8$~год сну. Чим така машина гірша за машину з гойдальними порогами, якщо одного разу вночі її розбудили через $4$~години?
\end{exercise}

\begin{exercise}[\lvlB]
\label{ex:20:4}
\emph{Період повільних гойдалок.} У моделі~\eqref{eq:updown} $F(x)=1/\bigl(1+e^{-(x-\theta)/k}\bigr)$, $w=5$, $I=2$, $\theta=2.5$, $k=0.25$, $b=5$, $\tau\ll\tau_a=0.3$~с. (а)~Знайдіть точки повороту $a_{\text{в}}$ і $a_{\text{н}}$ кривої $r=F(wr+I-a)$, де зникають робота й тиша. (б)~Узявши $r\approx1$ у роботі й $r\approx0$ у тиші, знайдіть період і частку роботи.
\end{exercise}

\begin{exercise}[\lvlB]
\label{ex:20:5}
\emph{Коли гойдалки вимикаються.} У моделі задачі~\ref{ex:20:4} нерухома точка лежить на перетині кривої $a=wr+I-F^{-1}(r)$ з прямою $a=br$. Коливання йдуть, поки перетин на середній гілці, між точками повороту. За яких $b$ це так? Як, змінюючи $b$, \nt{ACET} перемикає генератор на підсилювач?
\end{exercise}

\begin{exercise}[\lvlB]
\label{ex:20:6}
\emph{Моделювання: борг чи фаза.} У \texttt{sleep\_models.py} візьміть перше пробудження після $48$~год і тримайте машину без сну $D=24$, $32$, $40$, $48$, $64$~год (\texttt{stay\_awake}, \texttt{days=8}). Скільки триває відновлювальний сон за кожного $D$? Що визначає його тривалість?
\end{exercise}

\begin{exercise}[\lvlB]
\label{ex:20:7}
\emph{Моделювання: забування й перемішування.} Лінійний нейрон $y=\mathbf w\cdot\mathbf x$, $n=40$, учиться за правилом $\mathbf w\leftarrow\mathbf w-0.5\,(y-y^*)\,\mathbf x$. Задачі A і B --- по $10$ образів, $x_j\sim\mathcal N(0,1/n)$, $y^*=\pm1$. Яка середньоквадратична похибка на A після $200$ епох A, а потім $200$ епох B? А після $200$ епох A і B впереміш?
\end{exercise}

\begin{exercise}[\lvlC]
\label{ex:20:8}
\emph{Дослідження: перемішування чи перенормування?} Нейрон задачі~\ref{ex:20:7} з порогом; дві нічні процедури: (i)~усі ваги множаться на $0.82$; (ii)~старі образи повторюються впереміш із новими. Що кожна робить із відношеннями ваг і з відповідями нейрона? Як розрізнити гіпотези за розмірами синапсів після сну?
\end{exercise}
